\section{Cross-Shard Transaction Model}
\label{sec:cross-shard-model}

\subsection{System abstraction}
\label{subsec:system-abstraction}

We study a cross-shard transaction protocol in which a transfer moves value from
a source shard to a target shard while preserving a consistent terminal outcome.
Each shard maintains an independent UTXO-based ledger, and each cross-shard
transfer is represented by a uniquely identified protocol session.

A session coordinates source locking, cross-shard receipt creation and delivery,
destination preparation, and the final commit decision, while also handling
aborts, validation failures, insufficient quorum, timeouts, duplicate receipts,
and partially applied commits.

The implementation separates shard management from transaction-state mutation.
The shard manager maintains topology, validators, logical time, session
registration, and aggregate metrics, while the atomic-commit component validates
protocol preconditions and applies state changes. This separation allows the
transaction protocol to be analyzed independently from unrelated functionality
in the surrounding simulation platform.

Receipt authenticity is an abstract protocol precondition rather than a
cryptographic result. Quorum summarizes validator approval without modeling the
consensus mechanism that produces votes, and timeouts use logical rounds rather
than wall-clock deadlines.

\subsection{Cross-shard session state machine}
\label{subsec:state-machine}

Each transfer is represented by a finite-state session. The implementation
contains the following states:

\begin{itemize}
\item \texttt{CREATED}
\item \texttt{SOURCE\_\allowbreak LOCKED}
\item \texttt{RECEIPT\_\allowbreak CREATED}
\item \texttt{RECEIPT\_\allowbreak DELIVERED}
\item \texttt{DESTINATION\_\allowbreak PREPARED}
\item \texttt{COMMITTED}
\item \texttt{ABORTED}
\item \texttt{TIMED\_\allowbreak OUT}
\item \texttt{FAILED\_\allowbreak VALIDATION}
\end{itemize}

The normal successful execution follows

\begin{equation}
\begin{split}
\texttt{CREATED}
&\rightarrow \texttt{SOURCE\_\allowbreak LOCKED} \\
&\rightarrow \texttt{RECEIPT\_\allowbreak CREATED} \\
&\rightarrow \texttt{RECEIPT\_\allowbreak DELIVERED} \\
&\rightarrow \texttt{DESTINATION\_\allowbreak PREPARED} \\
&\rightarrow \texttt{COMMITTED}.
\end{split}
\label{eq:normal-commit-path}
\end{equation}

The states \texttt{COMMITTED}, \texttt{ABORTED}, \texttt{TIMED\_\allowbreak OUT},
and \texttt{FAILED\_\allowbreak VALIDATION} are terminal and have no outgoing
transitions.

Abort and validation-failure transitions may occur from nonterminal states. A
timeout may occur after source locking and before a terminal decision, but not
directly from \texttt{CREATED}, where no source resource has been locked.

The transition relation rejects commit after abort, commit after timeout,
creation-to-commit jumps, repeated terminal decisions, and receipt processing
for a different transfer identifier. Each accepted transition emits an
immutable event containing sequence number, logical time, transfer identifier,
action, previous state, next state, and diagnostic reason.

\subsection{Atomic commit and rollback}
\label{subsec:atomic-commit}

The executable commit operation has four conceptual stages:

\begin{enumerate}
\item preparation of the commit plan
\item validation of commit preconditions
\item application of ledger and session mutations
\item rollback if an application step fails
\end{enumerate}

Before modifying ledger state, the protocol constructs a commit plan and
captures a restorable snapshot containing the source UTXO and lock state,
receipt-consumption state, optional change output, destination output, and a
checkpoint of the complete protocol session.

During normal commit, the destination is prepared, the receipt consumed, the
source value removed, the source lock released, optional change created, the
destination output created, and the session moved to \texttt{COMMITTED}.

An application exception starts rollback from the captured snapshot. Successful
rollback restores the destination and change outputs, source UTXO, receipt
consumption, source lock, and session checkpoint. The journal records
\texttt{ROLLBACK\_COMPLETED}, and the propagated \texttt{ProtocolException}
reports \texttt{rolledBack=true}. If restoration throws, the journal records
\texttt{ROLLBACK\_FAILED}, the restoration error is attached to the application
error, and the propagated exception reports \texttt{rolledBack=false}.
That outcome gives no restoration guarantee. The session checkpoint can undo
an incomplete application without emitting a synthetic reverse transition;
the attempt journal retains intermediate stages in memory.

Deterministic failure-injection points around receipt consumption, source debit,
and destination credit make partially executed commits reproducible without
introducing nondeterministic faults.

\subsection{Concurrent formal abstraction}
\label{subsec:formal-abstraction}

% O1, formalización autocontenida del modelo histórico V1.
% Fuente ejecutable: specs/tla/CrossShardCommit.tla (paper1-q3-v1).
% No atribuir los 1272 trabajos históricos al modelo posterior V2.2.
\subsubsection{Executable semantics and the seven V1 property obligations}
\label{subsubsec:v1-executable-semantics}

The definitive \texttt{paper1-q3-v1} campaign evaluates the executable
\path{specs/tla/CrossShardCommit.tla} model, not the subsequent V2.2 model. Let $T$,
$S$, and $V$ denote the finite sets of transfers, shards, and validators.
The state vector is
\begin{equation}
\begin{split}
X=\langle &\mathit{status},\mathit{terminalStatus},\mathit{sourceShard},
\mathit{targetShard},\mathit{locked},\\
&\mathit{receiptOwner},\mathit{receiptUseCount},
\mathit{destinationCredit},\mathit{fundsReleased},\\
&\mathit{messages},\mathit{votes}\rangle.
\end{split}
\label{eq:editorial-v1-state}
\end{equation}

\texttt{Init} sets each transfer's status to \texttt{Pending}, terminal
status to \texttt{None}, lock/credit/release flags to false, receipt
uses to zero, and message/vote sets to empty. It chooses distinct source
and destination shards for each transfer. The executable transition
relation \texttt{Next} comprises \texttt{LockTransfer},
\texttt{ReleaseDelayedReceipt}, \texttt{ConsumeReceipt}, \texttt{CastVote},
\texttt{CommitTransfer}, \texttt{TimeoutTransfer}, and \texttt{Stutter},
with \texttt{CreditBeforeReceipt} and \texttt{CommitAfterAbort} guarded
by mutation modes and therefore disabled in the correct model. TLC
checks \(\mathit{Spec}=\mathit{Init}\land\Box[\mathit{Next}]_X\)
under the finite configurations supplied with the artifact.

In the correct model, \texttt{ConsumeReceipt} requires a locked transfer,
a non-delayed receipt message, zero previous uses, and no prior credit.
The same action increments the receipt-use counter and sets the Boolean
credit flag; this model does not expose an independently fallible
cross-shard destination-credit operation. \texttt{CommitTransfer}
requires a prepared transfer, destination credit, and at least the
configured quorum of votes. An enabled \texttt{TimeoutTransfer} in the
locked state marks abort and release in the same abstract transition.
These are substantial modeling restrictions, not implementation-level
proofs of distributed atomic commit.

The seven invariants checked in the V1 campaign are indexed below to
keep their full mathematical definitions readable. The index-to-source
correspondence is exact:
\begin{center}
\small
\begin{tabular}{ll}
$P_1$ & \texttt{NoReceiptReplay} \\
$P_2$ & \texttt{DestinationCreditRequiresValidReceipt} \\
$P_3$ & \texttt{DecisionConsistency} \\
$P_4$ & \texttt{EventuallyReleasedAfterTimeout} \\
$P_5$ & \texttt{QuorumRequired} \\
$P_6$ & \texttt{NoValueLossAtTermination} \\
$P_7$ & \texttt{TerminalStateIrreversibility}
\end{tabular}
\end{center}

Using the actual state variables of the executable V1 model, the seven predicates are
\begin{align}
P_1 &\equiv \forall t\in T:\mathit{receiptUseCount}[t]\leq1,\notag\\
P_2 &\equiv \forall t\in T:
  \mathit{destinationCredit}[t]\Rightarrow\notag\\
  &\qquad\bigl(\mathit{locked}[t]\land
  \mathit{receiptOwner}[t]=\mathit{targetShard}[t]\land\notag\\
  &\qquad\mathit{receiptUseCount}[t]=1\bigr),\notag\\
P_3 &\equiv \forall t\in T:
  \mathit{status}[t]=\texttt{Committed}\Rightarrow\notag\\
  &\qquad\neg\mathit{fundsReleased}[t],\notag\\
P_4 &\equiv \forall t\in T:
  \mathit{status}[t]=\texttt{Aborted}\Rightarrow\notag\\
  &\qquad\mathit{fundsReleased}[t],\notag\\
P_5 &\equiv \forall t\in T:
  \mathit{status}[t]=\texttt{Committed}\Rightarrow\notag\\
  &\qquad\lvert\mathit{votes}[t]\rvert\geq\mathit{Quorum},\notag\\
P_6 &\equiv \forall t\in T:\bigl(
  \mathit{status}[t]\in\{\texttt{Pending},\texttt{Locked},\texttt{Prepared}\}\notag\\
  &\qquad\lor (\mathit{status}[t]=\texttt{Committed}\land
     \mathit{destinationCredit}[t])\notag\\
  &\qquad\lor (\mathit{status}[t]=\texttt{Aborted}\land
     \mathit{fundsReleased}[t])\bigr),\notag\\
P_7 &\equiv \forall t\in T:\bigl(
     \mathit{terminalStatus}[t]=\texttt{None}\land\notag\\
  &\qquad\mathit{status}[t]\in
     \{\texttt{Pending},\texttt{Locked},\texttt{Prepared}\}\notag\\
  &\qquad\lor(\mathit{terminalStatus}[t]=\texttt{Committed}\land\notag\\
  &\qquad\mathit{status}[t]=\texttt{Committed})\notag\\
  &\qquad\lor(\mathit{terminalStatus}[t]=\texttt{Aborted}\land\notag\\
  &\qquad\mathit{status}[t]=\texttt{Aborted})\bigr).
\label{eq:editorial-v1-invariants}
\end{align}

These formulas use the actual state variables of the V1 TLA+ module. In
particular, $P_4$ is an invariant about the \texttt{Aborted} state rather
than a temporal eventuality. The formulas neither introduce new
verification outcomes nor retroactively alter the 1,272-task protocol.

Moreover,
several properties follow closely from action guards and simultaneous
state updates. Detecting a mutant that removes exactly such a guard
therefore does not independently establish implementation correctness,
non-vacuity of every property, or adequacy against broader defect classes.
The V1 state has a transfer-level terminal decision and Boolean
credit/release indicators; it cannot express two independently
irreversible decisions, one per shard, or failures between all concrete
ledger updates.

The later TLA+ V2.2 specification introduces distinct
\texttt{srcState}, \texttt{dstState}, and \texttt{decision} variables
and explicit \texttt{ApplySrc}/\texttt{ApplyDst} transitions. G2-P3
recorded three Java failure-injection traces with partial source debit
followed by an exact rollback, while V2.2 does not reverse an applied
commit debit. Thus neither one-to-one internal correspondence nor a
general Java--V2.2 refinement proof is claimed. The V2.2 source and
configurations support the semantic comparison with the Java witnesses.
The available V2.2 source copy lacks the original logs for the 52-case
matrix, so this manuscript reports no execution totals or verdicts from
that matrix. The Java G1--G2 observations are supplementary to the frozen
1,272-task V1 campaign.


Bounded TLA+ and Alloy models abstract the executable protocol to the
protocol-level state required for concurrent sessions and failure behavior. They
represent transfer status, source and target shards, resource locks, receipt
ownership and consumption, destination credit, released funds, messages, and
validator votes, while collapsing implementation states that do not alter the
protocol-level decision state.

\begingroup\raggedright
The Java-to-formal projection maps \texttt{CREATED} to \texttt{Pending};
\texttt{SOURCE\_\allowbreak LOCKED}, \texttt{RECEIPT\_\allowbreak CREATED}, and
\texttt{RECEIPT\_\allowbreak DELIVERED} to \texttt{Locked};
\texttt{DESTINATION\_\allowbreak PREPARED} to \texttt{Prepared};
\texttt{COMMITTED} to \texttt{Committed}; and \texttt{ABORTED},
\texttt{TIMED\_\allowbreak OUT}, and \texttt{FAILED\_\allowbreak VALIDATION} to
\texttt{Aborted}.
\par\endgroup

TLA+ additionally stores the first abstract terminal decision in
\texttt{terminalStatus}. Alloy instead represents bounded executions as an
ordered sequence of states and checks terminal-state irreversibility over
consecutive states. The two models therefore encode terminal consistency
differently while targeting the same protocol concern.

Concurrent configurations vary shards and active transfers and include
duplicate receipts, delayed messages, timeout races, and quorum behavior. They
are finite by construction, so conclusions are bounded by the declared scopes
and constants.

TLA+ and Alloy are complementary representations of the same protocol concerns,
not equivalent state spaces or interchangeable cost models. Absolute
execution-time differences are therefore not used to infer tool superiority.

\subsection{Evaluated protocol properties}
\label{subsec:properties}

Seven protocol properties define the main bounded verification target.

\begin{description}

\item[] \texttt{NoReceipt\allowbreak Replay}.
A cross-shard receipt must not be consumed more than once.

\item[] \texttt{DestinationCreditRequires\allowbreak ValidReceipt}.
Destination-side credit requires a receipt accepted as valid by the protocol.

\item[] \texttt{Decision\allowbreak Consistency}.
A committed transfer must not simultaneously be represented as having released
its source funds through an abort-like terminal outcome.

\item[] \texttt{NoValueLossAt\allowbreak Termination}.
At a terminal state, transferred value must either have been committed to the
destination or remain recoverable on the source side.

\item[] \texttt{TerminalState\allowbreak Irreversibility}.
Once the protocol reaches its first terminal decision, a later execution step
must not replace it with an incompatible terminal outcome.

\item[] \texttt{EventuallyReleasedAfter\allowbreak Timeout}.
An abstract transfer represented as \texttt{Aborted} must have released its
source funds. This property captures the release obligation associated with the
modeled timeout and non-commit terminal behavior.
Despite its historical name, it is a state invariant, essentially
\texttt{Aborted} $\Rightarrow$ \texttt{fundsReleased}, not a general temporal
liveness property. The study does not establish fairness, eventual decision
or commit, or freedom from permanent locking for arbitrary executions.

\item[] \texttt{Quorum\allowbreak Required}.
A commit requires the relevant validator quorum. TLA+ parameterizes this
threshold through \texttt{Quorum}, whereas the evaluated Alloy model requires at
least two validator votes.

\end{description}

Alloy checks assertions with the same seven names. Its replay assertion
excludes \texttt{ManyUses} for each receipt atom. Its credit assertion
requires a consumed receipt owned by the transfer, using the relation from
receipts to transfers. Its two terminal-state implications require committed
credit and aborted release. Its remaining assertions forbid committed release,
require two votes at commit, require release when aborted, and preserve
terminal status between adjacent states. In both V1 models these conditions
refer to abstract indicators. The mutation experiments target selected
assertions or invariants within these property families.

\subsection{Scope of the model}
\label{subsec:model-scope}

Receipt authenticity remains abstract, validator quorum is an approval
condition rather than a verified consensus protocol, and time is represented by
logical rounds. Network behavior is scenario-bounded: delay and duplicate
receipts are represented in evaluated configurations, while concrete network
events not structurally represented by the abstraction are projected as
stuttering steps during trace conformance.

Both formal representations operate over finite configurations. Absence of a
counterexample therefore means only that no violation was found for the
evaluated model, property, configuration, and bound, not an unbounded proof of
protocol correctness.

The Java implementation is related to the formal model through bounded trace
conformance, not a general refinement proof or behavioral equivalence.

% Interpretación formal: no modifica ni reejecuta los modelos congelados.
\subsection{Claim-level formalization and observation boundaries}
\label{subsec:claim-formalization}

Table~\ref{tab:three-levels} describes the guards and effects of the
unmutated systems. UTXO denotes an unspent transaction output. The two V1
models use abstract credit and release indicators; V2.2 tracks signed balance
deltas, and Java records individual outputs and owners.

\begin{table}[htbp]
\centering
\footnotesize
\caption{Guards, transition effects and observation boundaries of the evaluated
systems. The Java and V2.2 rows describe distinct execution contracts.}
\label{tab:three-levels}
\begin{tabular}{@{}p{0.12\linewidth}p{0.49\linewidth}p{0.32\linewidth}@{}}
\hline
Object & Guards, transitions and effects & Observation boundary \\
\hline
Java &
Begin validates the source UTXO, checks source quorum and obtains its lock. Commit requires a
nonterminal session, a delivered receipt (delivery can occur within the call),
a logical round within the deadline, destination approvals meeting quorum
and an unused receipt ID. The plan validates the source and output
preconditions and captures a snapshot. Application prepares the destination,
marks the receipt, debits and unlocks the input, creates change and destination
outputs, and records commit. Application failure attempts local restoration. &
Concrete UTXO amounts and owners, receipt IDs, locks and session state.
The in-memory journal exposes stages and distinguishes completed from failed
rollback. The protocol monitor and journal cover one JVM. \\
TLA+ V1 &
\texttt{LockTransfer} moves a pending transfer to locked and enqueues receipts.
\texttt{ConsumeReceipt} requires a locked transfer and an available, non-delayed
receipt with zero previous uses and no credit. It prepares the transfer,
counts the use and sets credit together. Distinct votes accrue after prepare.
Commit requires prepare, credit and quorum. Enabled timeout from locked
without credit aborts and releases in the same step. &
One abstract terminal status per transfer, with a first-decision marker.
Credit and release are flags. The replay projection groups concrete events
and permits stuttering. \\
Alloy V1 &
Adjacent ordered states obey the trace relation. Lock introduces receipt
messages. A queued receipt message for a locked transfer, with zero uses
for that receipt atom, enables prepare, one use and credit. Distinct votes
accrue after prepare. Commit requires prepare, credit and two votes.
Timeout from locked without credit aborts and releases. &
Finite ordered relational traces with per-receipt-atom use counts.
Adjacent-state stability represents terminal irreversibility. The scopes
and quorum differ from TLC. \\
TLA+ V2.2 &
\texttt{Lock} reserves an unspent, unowned input. \texttt{Prepare} consumes
an available receipt with no prior use or decision. \texttt{DecideCommit}
requires prepare, quorum and no decision. \texttt{DecideAbort} requires
a locked source and no decision. \texttt{ApplySrc} debits or releases a
locked source, and \texttt{ApplyDst} credits a prepared destination or
rejects an unprepared or prepared destination, according to that decision. &
Separate source, destination and decision states permit partial application.
Balances are signed deltas. The model has no transition that reverses an
applied debit to represent the observed Java rollback. \\
\hline
\end{tabular}
\end{table}

\paragraph{V1 obligations}
The \texttt{ConsumeReceipt} action changes the receipt-use counter and
\texttt{destinationCredit} in a single abstract transition; the two effects
cannot be independently interrupted at that level. The alias
\texttt{AtomicCommit} is exactly \texttt{DecisionConsistency}, whose
meaning is
\[
  \forall t:\ \texttt{status[t]}=\texttt{Committed}
  \Rightarrow \neg\texttt{fundsReleased[t]}.
\]
It is not a cross-node atomic-commit theorem. Similarly, the historical
\texttt{EventuallyReleasedAfterTimeout} formula is the state invariant
\(\forall t:\ \texttt{status[t]}=\texttt{Aborted}\Rightarrow
\texttt{fundsReleased[t]}\). Its name must not be read as a temporal
liveness result. Several of the seven evaluated safety properties follow
closely from their corresponding transition guards; removing precisely those
guards in an associated formal mutant tests targeted sensitivity, not
independent completeness of the property suite.

\paragraph{Where the seven V1 obligations obtain their support.}
Table~\ref{tab:v1-guards} makes the inductive dependencies explicit. A
property supported by a guard is meaningful within the transition system,
but a mutant that simply removes the same guard is a weak independent test
of specification adequacy. The original checker logs and reachability
witnesses are needed before claiming non-vacuous activation of each
antecedent.

\begin{table}[htbp]
\centering
\scriptsize
\caption{V1 target conditions and their direct transition dependencies;
these are source-level interpretations, not additional checker results.}
\label{tab:v1-guards}
\begin{tabular}{@{}p{0.38\linewidth}p{0.50\linewidth}@{}}
\hline
Evaluated property & Directly relevant modeled condition/action \\
\hline
\texttt{NoReceiptReplay} & \texttt{CanConsumeReceipt} prevents a second
increment of \texttt{receiptUseCount}. \\
\shortstack[l]{\texttt{DestinationCreditRequires}\\\texttt{ValidReceipt}} & \texttt{ConsumeReceipt}
sets owner, use count and destination-credit flag in one action. \\
\texttt{DecisionConsistency} & \texttt{CommitTransfer} and abort actions
constrain committed status versus release. \\
\texttt{NoValueLossAtTermination} & State disjunction accepts pending
states, committed credit, or aborted release, without amounts. \\
\texttt{TerminalStateIrreversibility} & Terminal-state transitions and the
first-decision marker preclude incompatible later status. \\
\texttt{EventuallyReleasedAfterTimeout} & \texttt{TimeoutTransfer} changes
\texttt{status} and \texttt{fundsReleased} together. \\
\texttt{QuorumRequired} & \texttt{CanCommit} requires enough votes before
\texttt{CommitTransfer}. \\
\hline
\end{tabular}
\end{table}

\paragraph{The distinct V2.2 abstraction.}
The V2.2 predicate named \texttt{Atomicity} excludes two contradictory
terminal-state combinations:
\[
\begin{split}
 \forall t:\quad&\neg(\texttt{srcState[t]}=\texttt{Debited}
                 \land\texttt{dstState[t]}=\texttt{Rejected}),\\
 &\neg(\texttt{srcState[t]}=\texttt{Released}
                 \land\texttt{dstState[t]}=\texttt{Credited}).
\end{split}
\]
It does not require both effects to occur in the same step. The separately
defined \texttt{ReachPartialApplication} query explores such intermediate
application states and is not a proof of global atomicity. V2.2 also defines
a numerical \texttt{Conservation} invariant for signed shard-balance changes,
conditional on termination of initiated transfers. The \texttt{bal} map is
initialized to zero: this is not a specification of all UTXO identities,
recipients and pre-existing balances. The preserved
\texttt{correct\_disjoint.cfg} refers to \texttt{Spec}, not
\texttt{LiveSpec}; the mere presence of \texttt{Termination} and
\texttt{NoPermanentLock} in the source does not establish either liveness
property under fairness.

\paragraph{Witnessed implementation--model mismatch.}
The G2-P3 diagnosis records real Java execution prefixes in which a source
UTXO is debited and an injected exception triggers successful in-memory
restoration of the checkpoint. At the V2.2 level, \texttt{ApplySrc} requires
a prior coordinator decision and the frozen model does not express that
same return to the earlier locked state following debit. Therefore a
step-preserving mapping of those internal events is \emph{not established}
under the examined abstraction. This observation neither disproves every
possible boundary-level refinement nor proves crash recovery. The journal
is in-memory, and the measured contract is strictly local to the JVM.

\paragraph{Non-vacuity obligations.}
For each conditional safety property, verifying its formula on a finite
state space is insufficient to show that its antecedent was reached.
Diagnostic witnesses for receipt consumption, commit, abort, timeout,
contention, and partial application must be associated with the exact model,
configuration, reachable state and checker result. The original V1
verification counts are not retroactively enlarged by later G2/P3 witnesses.
No proof of a new property is claimed in this paper without a corresponding
checker record and bound.
